% ------------------------------------------------------------
% §1  Opener: the Josephus problem as epistemological archetype
% ------------------------------------------------------------
\section{A pattern where none was promised}\label{sec:josephus}
% Save this section's hyperref anchor so citations appearing after the
% floated figure still record a section-level backref (not figure.1),
% keeping the bibliography's "Cited in §…" list free of duplicates.
\makeatletter\edef\josephussectionhref{\@currentHref}\makeatother

\emph{Concrete Mathematics} opens with a small massacre~\citep[\S 1.1]{GKP1994}.
Number $n$ people $1,2,\dots,n$ standing in a circle, and going around the
circle eliminate every second person until one survives. Which position
$J(n)$ survives? For $n=10$ the eliminations run
$2,4,6,8,10,3,7,1,9$, and the survivor is $J(10)=5$
(\cref{fig:josephus}). Computing small cases and staring at them reveals
that if $n = 2^m + \ell$ with $0 \le \ell < 2^m$, then
\begin{equation}\label{eq:josephus}
  J(2^m + \ell) \;=\; 2\ell + 1 ,
\end{equation}
which says something better in binary: \emph{$J(n)$ is the one-bit
cyclic left shift of $n$}. For $n = 10 = (1010)_2$ we get $J(n) =
(0101)_2 = 5$. An iterative algorithmic process ends up having a
solution involving a simple bit manipulation of the binary
representation of the input. (Hold on to one question as you read: what
happens if every \emph{third} person is eliminated instead? We will not
answer it until \cref{sec:josephus3}, where it becomes a parable about
the constants in \SHA{}'s own definition.)

\begin{figure}[!ht]
  \centering
  % One circle of 10, with the given people crossed out (cumulative):
  % #1 = list of position/elimination-order pairs to cross, #2 = sublabel
  \newcommand{\joscircle}[2]{%
    \begin{tikzpicture}[baseline=(current bounding box.center)]
      \foreach \i in {1,...,10} {
        \pgfmathsetmacro{\ang}{90 - (\i-1)*36}
        \node[circle, draw, minimum size=4.5mm, inner sep=0pt,
              font=\scriptsize] (p\i) at (\ang:1.35) {\i};
      }
      \foreach \i/\elim in {#1} {
        \pgfmathsetmacro{\ang}{90 - (\i-1)*36}
        \draw[gray] ($(p\i)+(-0.19,-0.19)$) -- ($(p\i)+(0.19,0.19)$);
        \node[gray, font=\tiny] at (\ang:1.8) {\elim};
      }
      \node[font=\footnotesize, align=center] at (0,-2.35) {#2};
    \end{tikzpicture}}%
  \joscircle{2/1, 4/2, 6/3, 8/4, 10/5}{first lap:\\$2,4,6,8,10$ go}%
  \,$\to$\,%
  \joscircle{2/1, 4/2, 6/3, 8/4, 10/5, 3/6, 7/7}{second lap:\\$3,7$ go}%
  \,$\to\cdots\to$\,%
  \begin{tikzpicture}[baseline=(current bounding box.center)]
    \foreach \i in {1,...,10} {
      \pgfmathsetmacro{\ang}{90 - (\i-1)*36}
      \node[circle, draw, minimum size=4.5mm, inner sep=0pt,
            font=\scriptsize] (p\i) at (\ang:1.35) {\i};
    }
    \foreach \i/\elim in {2/1, 4/2, 6/3, 8/4, 10/5, 3/6, 7/7, 1/8, 9/9} {
      \pgfmathsetmacro{\ang}{90 - (\i-1)*36}
      \draw[gray] ($(p\i)+(-0.19,-0.19)$) -- ($(p\i)+(0.19,0.19)$);
      \node[gray, font=\tiny] at (\ang:1.8) {\elim};
    }
    \node[circle, draw, very thick, fill=projfill,
          minimum size=4.5mm, inner sep=0pt, font=\scriptsize]
          at (90-4*36:1.35) {5};
    \node[font=\footnotesize, align=center] at (0,-2.35)
      {end: $5$ survives};
  \end{tikzpicture}
  \caption{Circular decimation for $n=10$, read left to right; the
  ellipsis elides the intermediate eliminations of the final lap. Small
  gray indices give the global elimination order; position $5$ survives,
  and $J(10) = 5 = (0101)_2$ is the one-bit cyclic shift of
  $n = 10 = (1010)_2$.}
  \label{fig:josephus}
\end{figure}
% Restore the section anchor so citations after this float record
% section-level backrefs (not figure.1), keeping "Cited in §…" deduped.
\makeatletter\let\@currentHref\josephussectionhref\makeatother

Dwell for a moment on what \emph{kind} of fact this is. The problem is
defined \emph{iteratively}: the survivor is specified as the outcome of
$n-1$ eliminations performed in order, and the obvious way to compute
$J(n)$ is to simulate them, one by one, around the circle. The closed
form abolishes the simulation. The answer is obtained by merely
\emph{rearranging the bits of the input}---one cyclic shift, no iteration
at all: a computation that appeared to require $n-1$ sequential steps
collapses to a single stroke. Nothing in the problem statement promised
that such a collapse existed; it was simply there, waiting to be noticed.

It will pay immediately to name this operation the way the \SHA{}
standard does~\citep{FIPS180-4}. For a $w$-bit word $x$ and $0 \le r < w$,
write
\[
  \mathrm{ROTL}^{r}(x), \qquad \mathrm{ROTR}^{r}(x), \qquad \mathrm{SHR}^{r}(x)
\]
for the cyclic left rotation, the cyclic right rotation, and the ordinary
(bit-discarding) right shift of $x$ by $r$ positions. In this notation the
Josephus law \cref{eq:josephus} is simply
\[
  J(n) \;=\; \mathrm{ROTL}^{1}(n),
\]
the rotation taken in the $(\lfloor \lg n \rfloor + 1)$-bit window
that $n$ occupies. The rotation operators are \emph{literal components
of \SHA{}}. Its round function stirs the state with fixed maps
$\Sigma_0, \Sigma_1, \sigma_0, \sigma_1$, each an exclusive-or of two
or three rotations of a $32$-bit word (\cref{sec:roundshape})---the
survivor map of a circle of $32$ bits, composed with itself in
threes. The same innocent operation that \emph{creates} the visible
structure in the Josephus problem is deployed in \SHA{} to
\emph{destroy} visible structure; and one of the very few theorems
known about \SHA{}'s internals is a statement about exactly these
xor-of-rotation maps (\cref{sec:rivest}).

\emph{Innocent discrete processes secretly respect structure}, and the
structure is usually found empirically---by computing small cases,
arranging them suggestively, and noticing. The recurrence came first; the
closed form \cref{eq:josephus} was extracted from data.

The Josephus problem is not a lone curiosity. A whole gallery of famous
iterative processes collapse the same way, each to a one-line function
of the binary digits of its input: the Tower of Hanoi, Gray codes, the
game of Nim, the parity of Pascal's triangle, even---in a sense that
must be stated carefully---the notorious $3x+1$ map.
\Cref{app:gallery} collects these specimens, in increasing order of
discomfort: three where the collapse is total, one (Kummer's carry
theorem) where the collapse names the central character of this whole
paper, and one where a collapsing change of coordinates provably
\emph{exists} and provably \emph{helps not at all}. The moral compounds
across the gallery: bit operations are not just where these collapses
land, they are a recurring \emph{language} in which discrete iteration,
when it is secretly simple, tends to become simple---and
(\cref{sec:statespace}) they are precisely the alphabet from which
\SHA{} is built.

This paper is about the mirror image of that phenomenon. The function
\SHA{}, which we will introduce slowly, is a discrete process on binary
strings that was \emph{engineered so that no analyst would ever have a
Josephus moment with it}: no change of representation, no clever
coordinate system, no generating function is supposed to reveal any
respected structure whatsoever---even though the function is built from a
few dozen lines of elementary operations, all of them individually as
transparent as the decimation rule. \SHA{} too is defined iteratively---a
fixed round transformation applied sixty-four times, itself chained block
after block---and the working conjecture is that, unlike the decimation,
its iteration admits no collapse. Since the entire document is an
extended commentary on that claim, it is worth displaying once, in the
form of a conjecture---informal, but falsifiable in every clause:

\begin{conjecture}[the no-collapse conjecture; informal]
\label{conj:nocollapse}
No change of representation tames \SHA{}: no bit-shift pattern, no
closed form, no algorithm of any kind computes it materially faster
than performing its sixty-four rounds; and no nontrivial property of
its outputs can be predicted materially more cheaply than by computing
them.
\end{conjecture}

\noindent
Nothing remotely like this is proved, and \cref{sec:worlds} explains
why a proof is beyond the current reach of mathematics. What is known
traces the conjecture's boundary. Inward, the few honest theorems about
the internals say which components, taken alone, are transparent
(\cref{thm:rivest} is the loveliest); the round-by-round records of
\cref{sec:siege} measure how far into the iteration analysis currently
penetrates. Outward, the theory of time-bounded Kolmogorov complexity
(\cref{sec:liupass}) is the formal home of ``no shortcut exists,'' and
turns the conjecture's vague ``materially faster'' into exact
statements. In between lies everything this paper does. Whether the
engineers succeeded is, remarkably, not a settled question but an
empirical one, now running as an uncontrolled worldwide experiment for
a quarter of a century. Making parts of that experiment controlled,
small, and measurable is the research program of this note.
